\chapter{Inspect prediction and local reconstruction}
\MARSChapterHeading{Inspect prediction and local reconstruction}
\label{revision-lab:predicting-pictures}

Read Chapter 26 with this lab. Import the public student archive
\texttt{035-h263-prediction-0.1.0-student.zip} through
\textbf{Import MARS Lab...}, then open
\textbf{Predict and reconstruct H.263 blocks} in Learn. The package contains six
complete experiments and one continuing independent function. Each experiment
isolates a numerical decision before the next lab joins those decisions into a
bitstream. No video file or external decoder is needed for these experiments.

Answer the predictions in Learn before starting the offered Debug activity.
The source pointer identifies the relevant instruction; the named stops below
identify where to inspect its result. Read vector lanes at the stated width and
check what each function returns in \texttt{v1}: a score, candidate count,
coefficient count and clipping count have different units.

The complete encoder remains in
\path{revision/labs/files/h263-encoder.asm}; its ordinary entry needs the raw
files and controls supplied in the next unit. The original \texttt{035-h263}
construction package supports writing that encoder yourself. Its grading harnesses
are separate callers, not fragments to concatenate into a worked experiment.

\section{Establish the window before measuring it}

Start with \path{worked/sad-windows.asm}. Its caller supplies displaced origins
and a stride of 176. The first prediction asks you to distinguish the containing
extent from the number of active samples.

For \texttt{h263\_sad16} and its MSA counterpart, the pointers already select
current and reference origins; the stride is the advance between rows. Calculate
$15s+16$ for the minimal containing extent and list which bytes within it are
actually read. At stride 176, the extent is 2656 but only 256 samples belong to
the window. A guard or padding byte inside the allocation is still outside the
active rows.

Compare scalar and vector results for all-zero, all-255 and asymmetric inputs.
The maximum score is 65280. Explain the width of each vector stage: unsigned byte
magnitudes, halfword pair sums accumulated across rows, then wider horizontal
reduction.
At \texttt{sad\_vector\_accumulated}, each of eight halfwords in
\texttt{w4} is 8160: two differences per row, sixteen rows, each difference 255.
Pairwise reduction then gives four words of 16320 and two doublewords of 32640,
whose sum is 65280. At \texttt{lesson\_sad\_compared}, both saved scores have
that value and status is zero. At \texttt{lesson\_sad\_same\_window}, equal
pointers with equal strides produce zero. A wide final register cannot repair a
carry already lost in a byte lane.
\begin{figure}[!ht]
  \centering
  \begin{tikzpicture}[alt={Three crops of one actual light-theme debugger photograph. A identifies sad-windows.asm paused at line 185 before horizontal reduction. B shows the complete written w4 vector, hexadecimal 1fe0 repeated eight times. C interprets that same vector as eight halfwords, each hexadecimal 1fe0, decimal 8160. Each lane holds two differences per row times sixteen rows times 255. The complete score after reduction is 65280.}]
    \node[anchor=north west,inner sep=0,font=\small\bfseries] (title) at (0,0) {Widen before accumulating the row differences};
    \node[anchor=north west,inner sep=0] (frame) at ([yshift=-3mm]title.south west) {\MARSRevisionCrop{revision/captures/h263-prediction/h263-native-sad-unit.png}{0.376\linewidth}{0.03828125}{0.22222222}{0.77812500}{0.75111111}};
    \node[anchor=east,font=\small\bfseries] at ([xshift=-1mm]frame.west) {A};
    \node[anchor=north west,inner sep=0] (vector) at ([yshift=-3mm]frame.south west) {\MARSRevisionCrop{revision/captures/h263-prediction/h263-native-sad-unit.png}{0.96\linewidth}{0.24843750}{0.09777778}{0.28281250}{0.87777778}};
    \node[anchor=east,font=\small\bfseries] at ([xshift=-1mm]vector.west) {B};
    \node[anchor=north west,inner sep=0] (lanes) at ([yshift=-3mm]vector.south west) {\MARSRevisionCrop{revision/captures/h263-prediction/h263-native-sad-unit.png}{0.96\linewidth}{0.26171875}{0.05555556}{0.26953125}{0.92111111}};
    \node[anchor=east,font=\small\bfseries] at ([xshift=-1mm]lanes.west) {C};
  \end{tikzpicture}
  \caption{The actual paused frame is line 185 (A); the editor's source pointer
    above the debugger remains on the earlier prediction question. The complete
    \texttt{w4} value (B) contains eight halfword sums (C). Each is
    $2\cdot16\cdot255=8160$, hexadecimal \texttt{1fe0}. Reading the same bits
    as bytes would split each sum into two unrelated pieces. Horizontal widening and
    reduction follow this pause and produce the complete score 65280.}
  \label{fig:revision-h263-native-sad-unit}
\end{figure}


Also compare the same input pointer under different strides. It is a read-only
alias, but not necessarily the same window.

Continue with \path{worked/motion-ties.asm}. At
\texttt{lesson\_motion\_interior}, the radius-two search counts 25 candidates; at $(0,0)$,
count nine. Derive the displacement intervals before computing reference
addresses. Predict how the count changes at the left edge away from a corner, then check
15 at \texttt{lesson\_motion\_edge}. The corner stop reports nine.
Flat and repeated patterns expose tie handling: centre wins a minimum tie,
otherwise increasing vertical and horizontal displacement decide. Reverse
traversal is valid when it explicitly preserves this ordering.

\section{Test fractional prediction before the transform}

In \path{worked/negative-motion.asm}, use full-luma displacement $(-1,-1)$ at macroblock origin $(16,32)$. The chroma
origin is initially $(8,16)$; its reference integer origin becomes $(7,15)$ with
both half phases set. Write the corresponding decomposition for components minus
two, zero and one. A negative odd component distinguishes floor from division
that truncates toward zero.

First inspect the smaller \path{worked/diagonal-prediction.asm} experiment.
The \texttt{h263\_predict8} output has sixty-four predicted bytes followed by
sixty-four signed residual halfwords. Input strides remain independent of those
tight output strides. Before inspecting a residual, select the correct native
byte order and signed halfword interpretation. A negative residual is not a large
positive brightness value.

For current sample 128 and reference neighbours 128,128,128,96, derive prediction
120 and residual eight. Then use neighbours zero, zero, zero, one to distinguish
one final diagonal rounding from two successive rounded averages. The result is
zero at \texttt{lesson\_single\_rounding}. In the first experiment, four
predictions are 120 and the others 128, so the residual magnitude sum is 32;
\texttt{lesson\_prediction\_values} also retains one residual eight for inspection.
Check the entire predicted block and signed residual array, not only their
aggregate SAD. Different wrong arrays can have the same magnitude sum.

With diagonal phase, the reference footprint is nine by nine while current reads
remain eight by eight. The six-block coordinator derives four Y origins and one
for each chroma plane, then publishes six records only after every child succeeds.
Its total magnitude is bounded by $6\cdot64\cdot255=97920$. Return to \path{worked/negative-motion.asm}: its six-block stop has predicted
bytes 128, residuals zero and total zero. The next call supplies a bad Cr stride.
At \texttt{lesson\_coordinator\_rejected}, status one and total zero mean
metadata rejection; the earlier 1152-byte result is retained. This call exercises
validation before child execution. It does not demonstrate recovery from a child
that fails after earlier children succeed.

\section{Keep four representations visible}

Open \path{worked/reconstruction-chain.asm}. Its first call transforms a
constant residual block of two. At \texttt{lesson\_forward\_dc}, DC is sixteen
and the returned count is 64 coefficients. The next call encodes an INTER block
with current value 130, prediction 128 and quantizer eight.
Record residual two, DC sixteen, selected level one, reconstructed DC 23, inverse
residual three and final sample 131 at
\texttt{lesson\_first\_reconstruction}. The returned one now counts nonzero
selected levels. Mark the transition at which information is
lost. The coefficient coder consumes level one; the inverse transform consumes
coefficient 23. Swapping those arrays changes the reconstructed picture even if
their storage width is the same.

Use an asymmetric block to check transform orientation. Compare scalar and MSA
modes against an independent reference, allowing the stated one-unit transform
accuracy. The finite examples do not replace the full inverse-accuracy gate.
That gate generates its inverse inputs independently of the submitted forward
transform and checks all six prescribed cohorts, the error measures and exact
zeros in both modes and byte orders.

For the block encoder's 448-byte result, offsets 0--127 contain selected signed
halfword levels, 128--255 reconstructed coefficients, 256--383 signed inverse
samples and 384--447 final bytes. A code at the first INTRA position has its
special DC interpretation. The corresponding INTER position is an ordinary level.
Do not apply one interpretation indiscriminately across both block types.

Use raw 132 and returned reconstruction 131 as the next prediction. Residual
one falls into the quantizer's zero region. At
\texttt{lesson\_second\_reconstruction}, the reconstruction stays 131 and the
nonzero-level count is zero. Compare with an encoder incorrectly predicting from
earlier raw 130. Separate quantization loss from prediction-state disagreement:
identical raw pictures do not prove an exact reconstructed-reference skip.

\section{Write the independent reconstruction function}

Before implementing, use \path{worked/signed-reconstruction.asm} to separate
sign extension, addition and clipping. At \texttt{lesson\_signed\_sum}, the
correction $-129$ added to prediction 128 gives $-1$, which clips to zero.
Reading the same stored halfword as unsigned gives 65407 and changes the sum.
The endpoint stop distinguishes 255, which is already representable, from 256,
which must clip and increment the count. At \texttt{lesson\_native\_stride},
\texttt{0x0000000100000008} has low word eight but is not an admitted stride.
Validate the complete native doubleword before narrowing it.

Keep one implementation in \path{independent/h263-reconstruct8.asm}.
Replace the deliberate unfinished \texttt{break} at
\texttt{h263\_reconstruct8}. Pass prediction bytes in \texttt{a0}, signed
inverse halfwords in \texttt{a1}, destination bytes in \texttt{a2} and the
settings address in \texttt{a3}. Settings contain three unsigned native
doublewords: prediction stride, inverse stride and output stride, in bytes.
The selected object is 24 bytes and requires eight-byte alignment.

Each image has eight rows and eight selected samples per row. Prediction and
output strides are independently eight through 4096. Inverse stride is even,
sixteen through 4096, and its pointer requires two-byte alignment. Byte pointers
need no additional alignment. All selected containing extents and settings must
have non-wrapping exclusive ends. Check full-width metadata before sample reads.
The mapped active memory is a caller precondition; do not probe arbitrary
unmapped addresses to discover whether they are valid.

Prediction, inverse and settings are disjoint. Output is disjoint from inverse
and settings. It may be exactly the prediction pointer when output and prediction
strides are identical; otherwise it is disjoint from prediction. Partial output
overlap is outside this interface. Exact aliasing does not imply zero work:
the prediction still needs its signed correction.

Success returns zero in \texttt{v0} and the number of strictly out-of-range sums
in \texttt{v1}, publishing exactly the clipped output bytes. Each inverse sample
must lie in $[-256,255]$. Metadata rejection returns one and zero before either
input is read. An out-of-range sample returns two and zero, retaining every
caller byte. Reading the admitted inputs during sample validation is permitted.
Keep padding, settings and inputs unchanged except for the selected successful
output when it aliases prediction. Do no console or file I/O. Preserve saved
general registers, frame pointer, global pointer, stack pointer and the supplied
return continuation. Entry stack alignment is sixteen bytes, with 8192 private
bytes available.

Choose either validation followed by reconstruction, or reconstruction into a
private sixty-four-byte staging block followed by publication. A late invalid
sample must leave the entire caller output unchanged under both disjoint and
exact-alias calls. If staging uses the stack, restore its pointer on every return.
Neither solution needs a second picture-sized allocation.

\section{Predict the demo, then challenge its boundaries}

Open the independent file and select its \texttt{main}. The caller fills all
prediction bytes with 128 and repeats correction row
$(-129,-128,-127,0,126,127,128,255)$. Before running, derive every output and
count the sums outside the representable sample interval. At
\texttt{recon\_demo\_return}, a completed function should return zero and 24.
The output row is $(0,0,1,128,254,255,255,255)$. Continuing copies status and
count into their named fields and exits without console output.

A sum of zero or 255 does not count as clipped. Negative corrections must remain
signed until after addition. Change the final inverse value to 256 and predict
the failure status and unchanged output. Repeat with output exactly aliasing the
prediction. A function that validates while directly storing can damage earlier
samples before discovering that last error.

Use the six offered \textbf{Check this program} activities on that same
source. They progress through sample arithmetic (60 case placements), exact aliases
(24), metadata (36), invalid samples (16), a two-call reuse case (2), and the full
suite (138). The final suite includes both byte orders. These are overlapping
profiles, not 276 distinct input scenarios. Open a hint when needed, edit and retry;
feedback applies to the current editor text. Returning to a worked experiment does
not erase the independent implementation. After an edit, renew execution checks
that referred to an earlier source revision.

For standalone feedback from the public repository, use:
\Needspace{8\baselineskip}
\MARSCodeLabel{Check reconstruction and retained state}
\begin{Verbatim}
java -jar "$CLI" grade \
  revision/labs/checks/h263-reconstruct8.kts \
  revision/labs/files/h263-reconstruct8.asm
\end{Verbatim}
The unfinished trap is expected to fail until implemented. Checks cover both byte
orders, unaligned byte pointers, minimum inverse alignment, independent padded
strides, full-width invalid metadata, exact prediction aliases, early and late
invalid samples, exact clipped counts and repeated calls. They compare results,
access windows and retained state rather than an instruction sequence.

The codec's numerical entries have separate package domains:
\Needspace{8\baselineskip}
\MARSCodeLabel{Check the complete local block pipeline}
\begin{Verbatim}
python3 packages/035-h263/grade.py \
  revision/labs/files/h263-encoder.asm --domain block \
  --cli "$CLI" --java "$JAVA"
\end{Verbatim}
Use \texttt{sad}, \texttt{motion}, \texttt{prediction}, \texttt{macroblock},
\texttt{transform} or \texttt{quantization} to isolate another completed stage.
The runner supplies separate grading storage. Reserve its documented label prefixes
and address arenas, and use the complete assignment's private stack requirements
for those entries. The independent function's smaller contract does not replace
them. Keep measured instruction work, host elapsed time and encoded quality as
different observations. The next lab connects these numerical decisions to syntax
and independent stream decoding.

\section{Extend the experiment to a whole picture}

To inspect the generator separately, optionally create raw input in a new directory:
\Needspace{8\baselineskip}
\MARSCodeLabel{Generate raw QCIF pictures and their description}
\begin{Verbatim}
python3 packages/035-h263/construction/h263-fixtures/fixture.py \
  --output-dir work/qcif
\end{Verbatim}
The description records dimensions, strides, plane offsets, clock steps and hashes.
Eight pictures occupy 304128 bytes. Locate picture one's Y sample at $(16,32)$
and the first byte of its Cb region. An RGB viewer cannot infer this planar layout
from a raw file alone. Keep display aspect ratio separate from stored sample
geometry. The following lab bundles the same generated input with its native
package and checks an actual restricted H.263 stream. These numerical experiments
do not establish interoperability.
